% !TEX program = xelatex
% 2026-2027学年第一学期杭州市高三教学质量检测 语文参考答案
\documentclass[11pt,a4paper]{ctexart}
\usepackage{ywexam}
\usepackage{multirow}

\examfooter{语文参考答案}

\begin{document}

\answertitle{2026-2027学年第一学期杭州市高三教学质量检测}{语文参考答案}

\sechead{一、阅读（75分）}

\subsechead{（一）阅读I（19分）}

\ans{1}{A【解析】B项，张冠李戴，“大多数时刻处于备战状态”的不是抗冷基因，而是“作物”，同时表述有误，原文是“在温暖季节空耗能量”；C项，因果倒置，“借助RGF多肽信号轴，研究者找到……”有误，原文是“倘若我们能找到这套防护系统的核心分子……最终解析出……RGF多肽信号轴”；D项，表述错误，高棚顶改为拱形顶主要是为了防止棚被风吹倒。}

\ans{2}{C【解析】A项，“缓慢降温则不会造成受精障碍及结实率下降”于文无据；B项，“只要选择了……就不会有损失”表述绝对；D项，果树打权是为了控制侧枝生长，以免其影响主干，冬天光合作用有限，所以要尽早打权，此项逻辑有误。}

\ans{3}{B【解析】“主要是为了表现科研的艰辛”错误。}

\ans{4}{参考答案：标题：激活RGF小肽后挽回损失／激活RGF小肽后损失减少／激活RGF小肽后产量提升（2分）。}
\subans{主体}{：番茄和水稻的图文（1分）；文字——番茄挽回了超过50\%的产量损失／番茄硕果累累，水稻挽回约18\%产量损失／水稻籽粒饱满（1分）。}

\ans{5}{参考答案：问题共性：\xq{1}问题都从材料中生发而来，拓展了材料中未能展开的部分。\xq{2}引起受众的阅读兴趣，引导受众学习科学家精神。\xq{3}语言表达清晰，指向明确。}
\scorenote{（评分标准：答出\xq{1}\xq{2}，每点2分；若只答出“有针对性”“有细节”“关注有趣的故事”等笼统的内容，总共最多给2分；若仅从语言单一角度作答，最多1分；其他答案言之成理也可酌情给分。）}
\subans{问题四}{：示例一“中国方案”的独特性和领先性体现在哪里？／示例二农民经验在RGF的发现过程中起到了怎样的作用？}
\scorenote{（评分标准：2分，其他答案言之成理也可酌情给分。）}

\subsechead{（二）阅读II（18分）}

\ans{6}{B【解析】奶奶的话语，既是坚守早年的约定（爷爷当年发誓“再也不去当打鼓匠”），也是源于多年生活的顾虑（鼓书在当地属于被人歧视的行当，有“鼓匠的崽崽闹丧场”的俗语，奶奶反对是为了维护安稳生活），但她不是“刻薄”“不近人情”之人。}

\ans{7}{示例：太阳从东边坡地爬上来，屋前的柳叶和桑叶在晨风中摇晃，大门对着东边，那道光直晃眼睛。世界上每一个角落的每一株草木都无法逃避，它们不能拒绝阳光，老罗也一样。}
\scorenote{（评分标准：\xq{1}所选景物要符合清晨乡村的特点，有适当细节描写，给2分；\xq{2}景物描写能关联人物情绪，体现出罗一木心愿落空后的落寞（或失落、压抑、被束缚）等，给2分；\xq{3}无语病、错别字，表达流畅，给1分。）}

\ans{8}{参考答案：亲近：}
\subans{\xq{1}}{本文开篇以孙子罗小民的视角切入，文中又穿插了爷孙的日常对话，如“真的要去唱？”“奶奶那怎么办？”等，使读者真切感受到祖孙之间的情感联结。}
\subans{\xq{2}}{本文细腻地描写了爷爷最后一次登上茶馆舞台的场景，使读者仿佛置身现场，感受到了热烈的氛围。}
\subans{疏离}{：}
\subans{\xq{1}}{本文主要以第三人称（全知视角）叙述，冷静平和，将罗一木置于被观察的位置，引导读者反思鼓书艺人在时代洪流中的挣扎。}
\subans{\xq{2}}{本文略写奶奶生病、离世等重大事件，叙事克制，没有刻意渲染情绪，拉开了读者与人物的情感距离。}
\scorenote{（评分标准：“亲近”或“疏离”任选一个维度，写出该维度下的两点即可，每点2分，共4分。）}

\ans{9}{参考示例一：不可删去。}
\subans{理由}{：\xq{1}本句是全篇的文眼，从丧妻之鼓升华为自祭之鼓，一个被压抑的鼓书艺人终于真正为没落的行当、落寞的自己发声，点出“最后的鼓王”的悲剧内核，照应了标题。\xq{2}该句叙事上以“直到第二天早上人们才明白”暗示结局、制造延宕，前文人们只把这次打鼓当作表演，结尾才点破真相，增强了小说的悲剧感。\xq{3}没有清晰交代前一天晚上发生的事情，给读者留下想象的空间。}
\subans{参考示例二}{：可以删去。}
\subans{理由}{：\xq{1}文末罗一木打鼓的故事已完结，无须再交代人们是否明白鼓声的内涵。\xq{2}文末“无不动容”的描写，表明人们沉醉于罗一木高超的技艺，与“最后的鼓王”形成呼应，使人物形象更为丰满。\xq{3}文末写到人们的反应就戛然而止，没有交代罗一木的命运，这种不确定性让结尾更具张力，给读者留下想象的空间。}
\scorenote{（评分标准：如果没有判断，扣1分；理由每点2分，共6分。）}

\subsechead{（三）阅读III（23分）}

\ans{10}{ADG【解析】对称性结构，“于……见……”。}

\ans{11}{D【解析】“洎牧以谗诛”的“诛”是（被）杀的意思，表被动，意思、用法均不同。}

\ans{12}{C【解析】“在原文中混入并宣扬了当时的一些邪说”错，邪说本就混杂在《左传》所采的旧史中，不是杜预作注时混入原文；杜预只是在注解时扩大其影响。}

\ans{13}{（1）只有郑樵略微有志于探求大义，可是（只会）辑录前人旧说的这类人，（却）吵吵嚷嚷地起来与他争辩。}
\scorenote{（评分标准：缀学、徒、嚣然、争，各1分，共4分。）}
\subans{（2）}{想要阐发《春秋》声讨乱臣贼子的宗旨，却不能做到公允持平，（人们）所说的（评价）过于严苛，指的就是这种情况。}
\scorenote{（评分标准：译出大意给1分，申、平、“所谓……也”，各1分，共4分。）}

\ans{14}{参考示例一：《六国论》：苏洵论六国灭亡，重在警示北宋朝廷切勿一味纳币苟安，契合章学诚提出的以“义”为核，不专以考究事与文的主张。作为史论，《六国论》不求全记史，只是援引了割地赂秦等关键史事。苏洵将六国之亡尽归于“赂秦”，主观倾向明显，评价偏严苛，未能做到宽严公允。}
\subans{参考示例二}{：《过秦论》：贾谊叙秦兴衰，重在总结秦国灭亡的教训，劝君主行仁政，契合章学诚提出的以“义”为核，不专以考究事与文的主张。作为史论，《过秦论》不求全记史，但直陈秦统一天下、陈涉起义等关键材料。贾谊为突出“仁义不施而攻守之势异也”，极力渲染秦的暴政，主观色彩浓厚，未能做到宽严公允。}
\scorenote{（评分标准：结合文本内容并体现材料中的原则，每点2分；只提炼原则，不结合文本，每点1分；共6分，超过字数扣1分。）}

\subsechead{（四）阅读IV（15分）}

\ans{15}{D【解析】“相信终有一日”错。}

\ans{16}{参考答案：\xq{1}“持节欲同苏武劲”，截取苏武持节不屈、坚守志节的侧面（1分），以劝勉友人，出使时保有苏武一般坚贞不屈的气节（2分）。}
\subans{\xq{2}}{“自笑餐毡典属国”，选取苏武餐毡受难、备尝艰辛的情节（1分），诗人借典自比，自嘲宦游旅途的奔波劳苦（2分）。}

\ans{17}{参考答案：（1）安能摧眉折腰事权贵，使我不得开心颜。}
\subans{（2）}{士不可以不弘毅，任重而道远。}
\subans{（3）}{春蚕到死丝方尽，蜡炬成灰泪始干。（李商隐《无题》）／落红不是无情物，化作春泥更护花。（龚自珍《己亥杂诗》）／捐躯赴国难，视死忽如归。（曹植《白马篇》）／鞠躬尽瘁，死而后已。（诸葛亮《后出师表》）}

\sechead{二、语言文字运用（15分）}

\ans{18}{C【解析】剑拔弩张：形容情势紧张，在此处最为合适。同室操戈：一家人操起刀枪自相残杀，比喻内部的争斗。枕戈待旦：枕着兵器以待天明，指时刻警惕，准备作战。如临深渊：如同站在深渊边缘，形容处事谨慎小心。}

\ans{19}{D【解析】A项，“也使文意不流畅”错；B项，“以偏概全”错；C项，“影响了整体语段的严谨性”错。}

\ans{20}{参考答案：（1）A：\xq{1}　B：闭合、合拢　C：\xq{3}\xq{4}（每点1分，共3分）}
\subans{（2）}{闭合、合拢——凑到一起——符合（每点1分，共3分）}

\ans{21}{参考示例：座右铭：先天下之忧而忧，岂因祸福避趋之。（1分）}
\subans{理由}{：前句写以天下之忧为先，后句写不因个人祸福而避责，逻辑通顺、衔接自然，表现了以天下为己任、不因个人境遇而轻易改变的坚定志向。（2分）}

\sechead{三、写作（60分）}

\ans{22}{参照高考作文评卷标准。}

\par\vspace{0.5em}
\begin{center}
\footnotesize
\setlength{\tabcolsep}{3pt}
\begin{tabular}{|c|c|p{2.8cm}|p{2.8cm}|p{2.8cm}|p{2.8cm}|}
\hline
\multicolumn{2}{|c|}{} & 一等（20$\sim$16分） & 二等（15$\sim$11分） & 三等（10$\sim$6分） & 四等（5$\sim$0分） \\
\hline
\multirow{2}{*}{基础等级} & 内容20分 & 符合题意\newline 中心突出\newline 内容充实\newline 思想健康\newline 感情真挚 & 符合题意\newline 主题明确\newline 内容较充实\newline 思想健康\newline 感情真实 & 基本符合题意\newline 中心基本明确\newline 内容单薄\newline 思想基本健康\newline 感情基本真实 & 偏离题意\newline 中心不明确\newline 内容不当\newline 思想不健康\newline 感情虚假 \\
\cline{2-6}
 & 表达20分 & 符合文体要求\newline 结构严谨\newline 语言流畅\newline 字迹工整 & 符合文体要求\newline 结构完整\newline 语言通顺\newline 字迹清楚 & 基本符合文体要求\newline 结构基本完整\newline 语言基本通顺\newline 字迹基本清楚 & 不符合文体要求\newline 结构混乱\newline 语言不通顺语病多\newline 字迹潦草难辨 \\
\hline
发展等级 & 特征20分 & 深刻\newline 丰富\newline 有文采\newline 有创意 & 较深刻\newline 较丰富\newline 较有文采\newline 较有创意 & 略显深刻\newline 略显丰富\newline 略显文采\newline 略显创意 & 个别语句有深意\newline 个别例子较好\newline 个别语句较精彩\newline 个别地方有深意 \\
\hline
\end{tabular}
\end{center}

\end{document}
